\documentclass[tikz,border=4pt]{standalone}
\usepackage{ctex}
\usepackage{amsmath}
\usepackage{pgfplots}
\pgfplotsset{compat=1.18}

% p10-03a: 二次函数性质图——增减区间（a>0，对称轴左侧递减，右侧递增）
% y = (x-2)^2 - 1，对称轴x=2，顶点(2,-1)
\begin{document}
\begin{tikzpicture}
\begin{axis}[
  width=8cm, height=7.5cm,
  axis lines=center,
  xlabel={$x$}, ylabel={$y$},
  every axis x label/.style={at={(axis cs:5.0,0)}, anchor=west},
  every axis y label/.style={at={(axis cs:0,5.0)}, anchor=south},
  xmin=-1.0, xmax=5.0,
  ymin=-1.8, ymax=5.0,
  xtick={-1,1,2,3,4},
  ytick={-1,1,2,3,4},
  tick style={color=black},
  ticklabel style={font=\small},
  clip=false,
]
  % y = (x-2)^2 - 1
  \addplot[blue, thick, domain=-0.7:4.7, samples=100] {(x-2)^2 - 1};
  \node[right, blue] at (axis cs:4.5, 3.5) {$y=(x-2)^2-1$};

  % 对称轴 x=2（虚线）
  \draw[gray, dashed, thick] (axis cs:2,-1.8) -- (axis cs:2,5.0);
  \node[above, gray, font=\small] at (axis cs:2,5.0) {$x=2$};

  % 顶点 (2,-1)
  \addplot[only marks, mark=*, mark size=3pt, red] coordinates {(2,-1)};
  \node[below, red] at (axis cs:2,-1) {$(2,\,-1)$};

  % 递减区间箭头（x<2，从右向左箭头在曲线上）
  \draw[->, blue!70, very thick] (axis cs:2.1,-1.0) -- (axis cs:0.5, 2.25) node[midway, left, blue, font=\small] {递减};

  % 递增区间箭头（x>2，从左向右箭头在曲线上）
  \draw[->, blue!70, very thick] (axis cs:1.9,-1.0) -- (axis cs:3.5, 2.25) node[midway, right, blue, font=\small] {递增};

  % 区间标注
  \node[below, font=\small, blue] at (axis cs:0.5, -0.15) {$x<2$};
  \node[below, font=\small, blue] at (axis cs:3.5, -0.15) {$x>2$};
\end{axis}
\end{tikzpicture}
\end{document}
